\begin{proposition}[C.3.1 (Proved): the circle action and local shear]
The effective circle subgroup $\rho=\R/\Z\subset\C/\Z$ of R6
acts freely outside its fixed double curve. That curve is a sphere,
and its two complex normal weights are $1,-1$. A monodromy
representative may be chosen $\rho$-equivariant. In marked real
coordinates on a regular torus fibre, its linear part is
$I-2\widehat\delta\otimes\gamma$; its affine correction is the
translation supplied by R2 and depends on the trivialization.
\end{proposition}

\begin{proof}
On a regular torus the circle acts by translation in the primitive
lattice direction $\widehat\delta$, and is therefore free.
At a multiple fibre, apply the invariant covector $\gamma$ to the
equation that a circle translation agrees with the $j$-th affine
deck generator to the power $a$. Since the deck translation has
$\gamma$-component $\pm1/m_j$, equality forces
$a\equiv0\pmod{m_j}$. Primitivity of $\widehat\delta$ then
forces the circle element to be the identity. This checks freeness
on the quotient using the actual twist vectors, including the root
centre; the argument is not just a statement about the product cover.

In a cusp toric chart, the relevant primitive edge directions are
$e_1,e_2,e_1-e_2$, with $e_2=\widehat\delta$.
Their nonzero determinants with $e_2$ have absolute value one.
Thus stabilizers off the edge in direction $e_2$ are trivial. Along
that edge the circle is fixed; the two normal characters in the
adjacent unimodular cones are $1$ and $-1$ because the base monomial
is invariant. The opposite-edge identification gives the double
curve with its two distinct triple-point endpoints, namely a
$\PP^1$. This verifies the toric step of R6 used here.

The circle acts on the entire family and commutes with the endpoint
isomorphism $F_2$. Averaging a Riemannian metric over this compact
circle makes the orthogonal horizontal distribution invariant.
Its parallel transport and hence the monodromy commute with $\rho$.
On a local marked torus chart,
$\Pi_{c-2}=\Pi_c(I-2\widehat\delta\otimes\gamma)$.
Inserting the identity complex-linear part and translation of R2
gives the asserted affine formula. A locally chosen lift of its
translation must not be multiplied by a real interpolation parameter
without checking its period cocycle; that operation need not descend
to the torus bundle.
\end{proof}

